\documentclass[book.tex]{subfiles}
\begin{document}
\chapter{Physical Theories for Fluids}

\label{chap:fluid_theory}
\noindent\rule{11cm}{0.4pt} \\
\textbf{Prerequisites:} \autoref{chap:qft_basics}  \\
\textbf{Difficulty Level:} ** \\
\noindent\rule{11cm}{0.4pt}
\vspace{5mm}

We introduce a simple formulation of fluid dynamics as a physical theory. This simple formulation will serve as a foundation for the more sophisticated theories of fluid dynamics we will present in subsequent chapters.

The state space of a fluid is given by the flow field of the fluid over some specified geometric region $D \subseteq \mathbb{R}^3$. Formally, this would be a vector field which we can define as type
\begin{align}
    D \to \mathbb{R}^3
\end{align}

Notice here that we assume that the fluid is incompressible and has constant density. In this basic model .

For this first simple model of fluid evolution, we will avoid dealing with the complexities of turbulence or other effects and assume that we are dealing with simple convective flow. %(\textcolor{red}{TODO: Define convective flow properly here}).

\begin{align}
    \frac{\partial u}{\partial t} + c \nabla \cdot u &= 0
\end{align}

We assume conservation of mass, momentum, and energy in this system. %\textcolor{red}{TODO: What are the symmetries of this equation? How can we write these conservation laws?}

We assume that we have the ability to draw observations of the underlying vector field $u$ subject to a noise model given by $\mathcal{N}(b, \sigma)$ where $b$ is a bias term and $\sigma$ is a noise term.

The model is parameterized by a set of initial conditions on the geometry
\begin{align}
    u_0: D \to \mathbb{R}^3
\end{align}
In the case of convective flow, this initial condition is often called the initial wave. The parameter $c$ is the initial speed of the wave.

\begin{center}
\begin{tabular}{|c|c|}
    \hline
    States & $\mathcal{S} = D \to \mathbb{R}^3 $ \\
    \hline
    Symmetries & $\mathcal{G} = \{\}$\\
    \hline
    Operators & $\mathcal{O} = \{ \frac{\partial u}{\partial t} + c \nabla \cdot u = 0\} $\\
    \hline
    Parameters & $\mathcal{W} = \{u_0: D \to \mathbb{R}^3, c\}$ \\
    \hline
    Constraints & $\mathcal{C} =  \{\}$ \\
    \hline
    Hyperparameters & $\mathcal{HP} = \{\}$  \\
    \hline
    Observables & $\mathcal{B} = \{u(x, t) + \mathcal{N}(b, \sigma) \}$  \\
    \hline
\end{tabular}
\label{phs:theory:static:particle}
\end{center}

We can convert this physical theory into a simple differentiable program 

\begin{lstlisting}[language=python]
class ConvectiveFluid(D $\subseteq$ $\mathbb{R}^3$, $u_0$: $D \to \mathbb{R}^3$)

  def simulate(t: $\mathbb{R}$):
    u : $D  \to \mathbb{R}^3$
    return solve($\partial$(u, t) + c $\nabla \cdot$(u) = 0, $u_0$, 0, t)
\end{lstlisting}

We can apply different methods to implement \lstinline{solve} such as finite differences, PINNs, neural operators and more. We will study certain concrete implementations in the coming chapters.

\end{document}